% Descripción del sistema para la SAT Competition (formato exigido: IEEE
% Proceedings, 1-2 páginas, PDF). Se redacta en inglés, el idioma de los
% proceedings. BORRADOR vivo: se completa a medida que cierran los
% experimentos (ADR-0006). Las marcas \pendiente{} deben desaparecer antes
% del envío.
\documentclass[conference]{IEEEtran}
\usepackage[utf8]{inputenc}
\usepackage[T1]{fontenc}
\usepackage[nopatch=eqnum]{microtype}
\usepackage{booktabs}
\usepackage{xcolor}
\usepackage[hidelinks]{hyperref}

% Macros compartidas por todos los documentos LaTeX del proyecto (ADR-0006).
\newcommand{\labesat}{\textsc{LabeSAT}}
\newcommand{\kissat}{\textsc{Kissat}}
\newcommand{\satsuma}{\textsc{satsuma}}
% Versión del documento: la fija CI con -usepretex, o se queda en «borrador».
\providecommand{\versiondoc}{borrador}

\newcommand{\pendiente}[1]{\textcolor{red}{[PENDING: #1]}}

\begin{document}

\title{\labesat{}: Kissat with Verifiable Symmetry Breaking\\
       \large (draft for SAT Competition 2027)}

\author{\IEEEauthorblockN{Abel Ponce}
\IEEEauthorblockA{\pendiente{affiliation}}}

\maketitle

\section{Introduction}

\labesat{} is a sequential CDCL solver for the Main Track. It is built from
\kissat{}~4.0.4~\cite{kissat} and, as a separate preprocessing program, uses
the symmetry-breaking tool \satsuma{}~\cite{satsuma}. \satsuma{}'s
substitution-redundancy (SR) proof prefix is continued by our \kissat{}
derivative, and the complete proof is checked against the original formula
with DSRtrim/Trestle~\cite{dsrtrim}.

\section{Solving pipeline}

A driver script first runs \texttt{satsuma fix} with time and size caps. It
then runs \labesat{}'s \kissat{} on the simplified formula with the new
application option \texttt{--append-proof}, which opens the proof file in
append mode. If \satsuma{} fails, exceeds its cap, or the input is too large,
the driver solves the original formula with \kissat{} alone and writes a pure
DRAT proof. The time spent in \satsuma{} is deducted from the solver's time
limit.

\satsuma{} is compiled without the optional clique component (cliquer, GPL),
so that the whole entry is MIT licensed. \pendiente{final decision D-005 and
its measured cost; see EXP-007}

\section{Changes to \kissat{}}

\begin{itemize}
  \item \texttt{--append-proof}: continue an existing (binary) proof file.
  \item A termination fix: the \emph{lucky} phases now return control to the
        solver's terminator (time/conflict limits).
  \item \pendiente{conditional activation of symmetry breaking (B3), if
        validated}
\end{itemize}

An adaptive stable/focused scheduler was implemented and evaluated, but it
showed no effect in a pre-registered experiment (factor $0.983\times$,
$p=0.69$). It is therefore \emph{disabled} in the submitted configuration.

\section{Mandatory statement}

\begin{itemize}
  \item \textbf{Code base}: \kissat{} 4.0.4 (upstream commit
        \texttt{8af8e56}); \satsuma{} commit \texttt{c6ad1b5} (unmodified);
        dejavu commit \texttt{4c275e9} (unmodified).
  \item \textbf{Amount of change}: 390 lines added and 10 removed in
        \kissat{}'s 38\,944 lines of C (1.0\,\%). Only 105 added lines are
        active in the submitted configuration; 285 belong to the disabled
        adaptive scheduler and its tracing. A 125-line driver script was also
        added. Figures produced by \texttt{scripts/declaracion\_ia.sh}.
        \pendiente{regenerate at submission}
  \item \textbf{Changes due to AI}: all of the added lines were written by
        an AI coding assistant (Claude Code) under the author's direction.
        The experimental methodology was also AI-assisted.
        \pendiente{final wording after the organizers answer (D-002)}
  \item \textbf{Heuristics tuned by AI}: none by automated search. The
        \satsuma{} time/size caps were fixed in a pre-registered experiment
        designed with AI assistance. \pendiente{confirm}
\end{itemize}

\section*{Acknowledgments}
We thank A.~Biere and colleagues for \kissat{}, M.~Anders for \satsuma{} and
dejavu, and C.~R.~Codel, J.~Avigad and M.~J.~H.~Heule for DSRtrim and
Trestle.

\begin{thebibliography}{9}
\bibitem{kissat} A.~Biere, T.~Faller, K.~Fazekas, M.~Fleury, N.~Froleyks and
  F.~Pollitt, ``CaDiCaL, Gimsatul, IsaSAT and Kissat entering the SAT
  Competition 2024,'' in \emph{Proc.\ SAT Competition 2024}, 2024, pp.~8--10.
\bibitem{satsuma} M.~Anders, \emph{satsuma}, \url{https://github.com/markusa4/satsuma}.
\bibitem{dsrtrim} C.~R.~Codel, J.~Avigad and M.~J.~H.~Heule, ``Verified
  substitution redundancy checking,'' SAT Competition 2026 proof checker
  documentation.
\end{thebibliography}

\end{document}
